\subsection{SZX diagram families and list instantiation}

We introduce the definition of a family of SZX diagrams $D: \N^k \to \diag$
as a function from $k$ integer \textit{parameters} to SZX diagrams.
We require the structure of the diagrams to be the same for all elements in the family,
parameters may only alter the wire tags and spider phases.
Partial application is allowed, we write $D(n)$ to fix the first parameter of $D$.

\begin{example}
    The following example describes a family of diagrams $D$ that applies z-rotations
    with angle $\sfrac{\pi}{n}$ on $n+1$ qubits.
    \[D := n \mapsto \tikzfig{szx/family-example}\]
\end{example}

Since instantiations of a family share the same structure,
we can compose them in parallel by merging the different values of wire tags and spider phases.
We introduce a shorthand for instantiating a family of diagrams on multiple values
and combining the resulting diagrams in parallel.
This definition is strictly more general than the \textit{thickening endofunctor}
presented by Carette et al.~\cite{carette_quantum_2021},
which replicates a concrete diagram in parallel.
A \textit{list instantiation} of a family of diagrams $D: \N^{k+1} \to \diag$ over
a list $N$ of integers is written as $(D(n), n \in N)$.
This results in a family with one fewer parameter, $(D(n), n \in N): \N^k \to \diag$. 
We graphically depict a list instantiation as a dashed box in a diagram, as follows.
\[\tikzfig{szx/list-instantiation-box} := \tikzfig{szx/list-instantiation}\]

The definition of the list instantiation operator is given recursively
on the construction of $D$ in Table~\ref{tab:list-instantiation}.
On the diagram wires we use $v(N)$ to denote the wire cardinality $\sum_{n \in N} v(n)$,
$\overrightarrow\alpha(N)$ for the concatenation of phase vectors
$\overrightarrow\alpha(n_1) :: \dots :: \overrightarrow\alpha(n_m)$,
and $\sigma(N)$ for the composition of permutations $\bigotimes_{n \in N} \sigma(n)$.
In general, a permutation arrow $\sigma(N,v,w)$ instantiated in concrete values
can be replaced by a reordering of wires between two gather gates
using the rewrite rule $\bf{(p)}$.

With these definitions, we can formalise the parallel ($\otimes$) and sequential ($\circ$) composition of families of diagrams.

\begin{definition}
    Given diagram families $D_i: \N^{k+1} \to \diag$, $N = [n_1, \dots, n_m] \in \N^m$. We define the parallel composition as:
    \[
        ((D_1 \otimes D_2)(n), n \in N) := (D_1(n), n \in N) \otimes (D_2(n), n \in N)
    \]
    In the same vein, we define the sequential composition as:
    \[
        ((D_2 \circ D_1)(n), n \in N) := (D_2(n), n \in N) \circ (D_1(n), n \in N) 
        %
    \]
\end{definition}

\begin{table}
    \[
        \tikzfig{szx/list-instantiation/wire}
        :=
        \tikzfig{szx/list-instantiation/wire-flat}
        %
        \qquad
        %
        \tikzfig{szx/list-instantiation/ground}
        :=
        \tikzfig{szx/list-instantiation/ground-flat}
    \]
    \[
        \tikzfig{szx/list-instantiation/hadam}
        :=
        \tikzfig{szx/list-instantiation/hadam-flat}
        %
    \]
    \[  %
        \tikzfig{szx/list-instantiation/arrow}
        :=
        \tikzfig{szx/list-instantiation/arrow-flat}
    \]
    \[
        \tikzfig{szx/list-instantiation/spider-Z}
        :=
        \tikzfig{szx/list-instantiation/spider-Z-flat}
        %
    \]
    \[
        %
        \tikzfig{szx/list-instantiation/spider-X}
        :=
        \tikzfig{szx/list-instantiation/spider-X-flat}
    \]
    \[
        \tikzfig{szx/list-instantiation/gather}
        :=
        \tikzfig{szx/list-instantiation/gather-flat}
    \]
    {\small Where $\sigma(N, v, w) \in \Ftwo^{v(N)+w(N) \times v(N)+w(N)}$ is the permutation defined as the matrix}
    \[
        \sigma(N, v, w) = \begin{pmatrix}\sigma_f^N \vert \sigma_g^N \end{pmatrix}, \quad
        \sigma_f^{N} \in \Ftwo^{v(N)+w(N) \times v(N)},
        \ \sigma_g^{N} \in \Ftwo^{v(N)+w(N) \times w(N)}
    \]
    \[
    \begin{array}{ll}
        \sigma_f^{[\,]} = Id_0
        &
        \sigma_f^{n::N'} =
            \begin{pmatrix}
                Id_{v(n)} & 0 \\
                0 & 0 \\
                0 & \sigma_f^{N'}
            \end{pmatrix}
        \\[10pt]
        \sigma_g^{[\,]} = Id_0
        &
        \sigma_g^{n::N'} =
            \begin{pmatrix}
                0 & 0 \\
                Id_{w(n)} & 0 \\
                0 & \sigma_g^{N'}
            \end{pmatrix}
    \end{array}
    \]
    \caption{Definition of the list instantiation operator.}%
    \label{tab:list-instantiation}
\end{table}

\begin{lemma}%
    \label{lem:list-instantiation}
    For any diagram family $D$, $n_0 : \N$, $N : \N^k$,
    \[
        \tikzfig{szx/list-instantiation-append}
        \;=\;
        \tikzfig{szx/list-instantiation-append-split}
    \]
\end{lemma}

\begin{lemma}%
    \label{lem:list-init-zero}
    A diagram family initialized with the empty list corresponds to the empty map.
    For any diagram family $D$,
    \[
        \tikzfig{szx/list-instantiation-empty}
        \;=\;
        \tikzfig{szx/list-instantiation-empty-flat}
    \]
\end{lemma}

\begin{lemma}%
    \label{lem:list-instantiation-linear}
    The list instantiation procedure on an $n$-node diagram family adds
    $\bigo(n)$ nodes to the original diagram.
\end{lemma}
